% !TEX root = MM2025.tex


%%%%%%%%%%%%%%%%%%%%%%%%%%%%%%%%%%%%%%%%%%%%%%%%%%
\clearpage
\section{B. {\sectionthree} Appendix\label{app:gaps}}
\renewcommand*{\thepage}{\Alph{section}\arabic{page}}
\setcounter{page}{1}
%%%%%%%%%%%%%%%%%%%%%%%%%%%%%%%%%%%%%%%%%%%%%%%%%%


%%%%%%%%%%%%%%%%%%%%%%%%%%%%%%%%%%%%%%%%%%%%%%%%%%
\subsection{Detailed AKM Estimation Results\label{app:akm}}
%%%%%%%%%%%%%%%%%%%%%%%%%%%%%%%%%%%%%%%%%%%%%%%%%%

% HOURS EFFECTS:
\begin{figure}[htbp!]
    %
    \includegraphics[width=0.66\columnwidth]{figures/figureB1.eps}%
    %
    \caption{\label{fig:akm_hours}Predicted AKM contractual work hours FEs, by gender}
    %
    \begin{figurenotes}
        %
        \footnotesize
        %
        Figure shows predicted AKM contractual work hours FEs separately for men and women based on estimating earnings equation \eqref{eq:akm}. The omitted category is 1 hour, for which the FE value is normalized to 0.
    \end{figurenotes}
    \begin{figurenotes}[Source]
        \sourcedata%
    \end{figurenotes}
\end{figure}

% OCCUPATION EFFECTS:
\begin{figure}[htbp!]
    %
    \includegraphics[width=0.66\columnwidth]{figures/figureB2.eps}%
    %
    \caption{\label{fig:akm_occ}Predicted AKM occupation FEs, by gender}
    %
    \begin{figurenotes}
        %
        \footnotesize
        %
        This figure shows the predicted AKM occupation FEs separately for men and women based on estimating the earnings equation \eqref{eq:akm}. The fixed effects of both genders are sorted by the mean FEs of men's FE quantiles. %
    \end{figurenotes}
    \begin{figurenotes}[Source]
        \sourcedata%
    \end{figurenotes}
\end{figure}

% ACTUAL-EXPERIENCE EFFECTS:
\begin{figure}[htbp!]
    %
    \includegraphics[width=0.66\columnwidth]{figures/figureB3.eps}
    %
    \caption{Predicted AKM actual-experience FEs, by gender\label{fig:akm_exp_act}}
    %
    \begin{figurenotes}
        %
        \footnotesize
        %
        This figure shows predicted AKM actual-experience FEs separately for men and women based on estimating the earnings equation \eqref{eq:akm}. %
    \end{figurenotes}
    \begin{figurenotes}[Source]
        \sourcedata%
    \end{figurenotes}
\end{figure}

% TENURE EFFECTS:
\begin{figure}[htbp!]
    %
    \includegraphics[width=0.66\columnwidth]{figures/figureB4.eps}
    %
    \caption{Predicted AKM tenure FEs, by gender\label{fig:akm_tenure}}
    %
    \begin{figurenotes}
        %
        \footnotesize
        %
        This figure shows predicted AKM tenure FEs separately for men and women based on estimating earnings equation \eqref{eq:akm}. %
    \end{figurenotes}
    \begin{figurenotes}[Source]
        \sourcedata%
    \end{figurenotes}
\end{figure}

% EDUCATION-YEAR EFFECTS:
\begin{figure}[htbp!]
    %
    \subfloat[Men]{\includegraphics[width=0.50\columnwidth]{figures/figureB5a.eps}}%
    \subfloat[Women]{\includegraphics[width=0.50\columnwidth]{figures/figureB5b.eps}}
    %
    \caption{Predicted AKM education-year FEs, by gender\label{fig:akm_edu_year}}
    %
    \begin{figurenotes}
        %
        \footnotesize
        %
        Figure shows predicted AKM education-year FEs separately for men and women based on estimating earnings equation \eqref{eq:akm}. Note that the declining pattern for both genders and all education categories is due to measuring earnings in multiples of the prevailing minimum wage, which increased over this period---see \citet{EngbomMoser2022a} for details. %
    \end{figurenotes}
    \begin{figurenotes}[Source]
        \sourcedata%
    \end{figurenotes}
\end{figure}

% EDUCATION-AGE EFFECTS:
\begin{figure}[htbp!]
    %
    \subfloat[Men]{\includegraphics[width=0.50\columnwidth]{figures/figureB6a.eps}}%
    \subfloat[Women]{\includegraphics[width=0.50\columnwidth]{figures/figureB6b.eps}}
    %
    \caption{Predicted AKM education-age FEs, by gender\label{fig:akm_edu_age}}
    %
    \begin{figurenotes}
        %
        \footnotesize
        %
        Figure shows predicted AKM education-age FEs separately for men and women based on estimating earnings equation \eqref{eq:akm}. Age-pay profiles for all education groups are constrained to be constant from age 45 to age 49 and unconstrained otherwise. %
    \end{figurenotes}
    \begin{figurenotes}[Source]
        \sourcedata%
    \end{figurenotes}
\end{figure}
%


%%%%%%%%%%%%%%%%%%%%%%%%%%%%%%%%%%%%%%%%%%%%%%%%%%
\clearpage
\subsection{Further Details on Between vs. Within-Employer Pay Differences\label{app:oaxaca_blinder}}
%%%%%%%%%%%%%%%%%%%%%%%%%%%%%%%%%%%%%%%%%%%%%%%%%%

Here, we present two alternative Kitagawa-Oaxaca-Blinder decompositions of the gender pay gap, only the first of which is presented in the main text. The overall gender gap in employer FEs can be written as
%
\begin{small}
    \begin{align}
        %
        &\mathbb{E}\left[\left.\psi_{MJ(i,t)}\right| G(i) = M\right] - \mathbb{E}\left[\left.\psi_{FJ(i,t)}\right| G(i) = F\right] \label{eq:gender_pay_gap} \\
        %
        &= %
        \underbrace{\left(\mathbb{E}\left[\left.\psi_{MJ(i,t)}\right| G(i) = M\right] - \mathbb{E}\left[\left.\psi_{MJ(i,t)}\right| G(i) = F\right]\right)}_{\text{between-employer gap}}%
        %
        + %
        %
        \underbrace{\mathbb{E}\left[\left.\psi_{MJ(i,t)} - \psi_{FJ(i,t)}\right| G(i) = F\right]}_{\text{within-employer gap}} \label{eq:oaxaca_blinder_A} %
        %
        \\
        %
        &= %
        \underbrace{\left(\mathbb{E}\left[\left.\psi_{FJ(i,t)}\right| G(i) = M\right] - \mathbb{E}\left[\left.\psi_{FJ(i,t)}\right| G(i) = F\right]\right)}_{\text{between-employer gap}}%
        %
        + %
        %
        \underbrace{\mathbb{E}\left[\left.\psi_{MJ(i,t)} - \psi_{FJ(i,t)}\right| G(i) = M\right]}_{\text{within-employer gap}} \label{eq:oaxaca_blinder_B} %
        %
        .
        %
    \end{align}
\end{small}
%

Table \ref{tab:oaxaca_blinder} shows the results from the two alternative decompositions of the gender log pay gap corresponding to equations \eqref{eq:oaxaca_blinder_A} and \eqref{eq:oaxaca_blinder_B}.

%
\begin{table}[htbp!]
    \caption{Alternative Kitagawa-Oaxaca-Blinder decompositions of the gender log pay gap}
    \label{tab:oaxaca_blinder}
    %
    \resizebox{\textwidth}{!}{%
    \input{tables/tableB1.tex}
    }
    %
    \begin{tablenotes}
        %
        \footnotesize
        %
        This table shows results from the Kitagawa-Oaxaca-Blinder decomposition of the overall gender log pay gap into a between-employer gap (i.e., pay-policy component) and a within-employer gap (i.e., sorting component). Decomposition 1 corresponds to equation \eqref{eq:oaxaca_blinder_A} and uses men's employer FEs for computing the between-employer component. Decomposition 2 corresponds to equation \eqref{eq:oaxaca_blinder_B} and uses women's employer FEs for computing the between-employer component. %
    \end{tablenotes}
    \begin{tablenotes}[Source]
        \sourcedata%
    \end{tablenotes}
\end{table}
%

To illustrate these two decompositions, Figure \ref{fig:oaxaca_blinder_components} shows the distributions of gender-specific employer FEs underlying the individual terms in equations \eqref{eq:oaxaca_blinder_A} and \eqref{eq:oaxaca_blinder_B}.

%
\begin{figure}[htbp!]
  %
  \subfloat[Decomposition 1: Between-gap, men's FEs]{\includegraphics[width=0.50\columnwidth]{figures/figureB7a.eps}\label{fig:oaxaca_blinder_components_A}}%
  \subfloat[Decomposition 1: Within-gap, women's weights]{\includegraphics[width=0.50\columnwidth]{figures/figureB7b.eps}\label{fig:oaxaca_blinder_components_B}} \\
  %
  \subfloat[Decomposition 2: Between-gap, women's FEs]{\includegraphics[width=0.50\columnwidth]{figures/figureB7c.eps}\label{fig:oaxaca_blinder_components_C}}%
  \subfloat[Decomposition 2: Within-gap, men's weights]{\includegraphics[width=0.50\columnwidth]{figures/figureB7d.eps}\label{fig:oaxaca_blinder_components_D}}
  %
  \caption{Components of Kitagawa-Oaxaca-Blinder decompositions\label{fig:oaxaca_blinder_components}}
  %
  \begin{figurenotes}
        %
        \footnotesize
        %
        This figure shows the pay distributions underlying the Kitagawa-Oaxaca-Blinder decompositions---specifically, %
        %
        the between-gap using men's FEs (Panel \subref{fig:oaxaca_blinder_components_A}) %
        and using women's FEs (Panel \subref{fig:oaxaca_blinder_components_C}), %
        %
        as well as %
        %
        the within-gap using women's weights (Panel \subref{fig:oaxaca_blinder_components_B}) %
        and using men's weights (Panel \subref{fig:oaxaca_blinder_components_D}). %
        %
        Decomposition 1 (Panels \subref{fig:oaxaca_blinder_components_A} and \subref{fig:oaxaca_blinder_components_B}) %
        and decomposition 2 (Panels \subref{fig:oaxaca_blinder_components_C} and \subref{fig:oaxaca_blinder_components_D}) %
        %
        correspond to equations \eqref{eq:oaxaca_blinder_A} and \eqref{eq:oaxaca_blinder_B}, respectively. The dashed vertical line shows the mean of the distribution. %
    \end{figurenotes}
    \begin{figurenotes}[Source]
        \sourcedata%
    \end{figurenotes}
\end{figure}
%
